\documentclass[12pt]{article}
\usepackage[utf8]{inputenc}
\usepackage{graphicx}
\usepackage{enumerate}
\newcommand{\pt}[1]{\textbf{(#1 pts)}}
\begin{document}
\includegraphics[scale=.25]{LogoMate.eps}

\vspace{-3ex}
$\phantom{x}$\hfill MAT003 Cálculo II\\
$\phantom{x}$\hfill Laboratorio 1, II Ciclo 2025\\[1ex]
Facultad de Ciencias Exactas y Naturales\hfill Duración: 2 horas\\
Lunes  1 de setiembre de 2025\\
Puntaje total: 4 puntos
\vspace{-1ex}
\hrulefill
\hfill v2 - SOLUCIONES

\section*{Pregunta 1 \pt{4}}
Calcule las siguientes medidas para la variable Visitas: D5 y D2.

\paragraph{Solución y retroalimentación:} El decil 5 (D5) de la variable Visitas corresponde al percentil 50. Para determinarlo se consideran las 100 observaciones ordenadas de menor a mayor y se obtiene el valor asociado con el 50 \% de la distribución. El resultado es 5,00 visitas. Esto indica que aproximadamente el 50 \% de las observaciones presentan valores de visitas menores o iguales a 5,00 visitas.

El decil 2 (D2) de la variable Visitas corresponde al percentil 20. Para determinarlo se consideran las 100 observaciones ordenadas de menor a mayor y se obtiene el valor asociado con el 20 \% de la distribución. El resultado es 1,00 visitas. Esto indica que aproximadamente el 20 \% de las observaciones presentan valores de visitas menores o iguales a 1,00 visitas.
\end{document}